\chapter{Wprowadzenie i analiza problemu}

\section{Opis analizowanego problemu}

Hotel przechowuje wiele informacji, które są ze sobą powiązane. Pracownik recepcji musi znać dane gości, dostępność pokoi, terminy rezerwacji, płatności oraz usługi dodatkowe. Jeżeli takie informacje są prowadzone tylko w notatkach albo arkuszu kalkulacyjnym, łatwo o pomyłkę, na przykład podwójną rezerwację tego samego pokoju w tym samym terminie.

Analizowany problem dotyczy podstawowego zarządzania hotelem w aplikacji desktopowej. Projekt nie jest pełnym systemem hotelowym klasy komercyjnej. Nie zawiera modułu płatności online, integracji z portalami rezerwacyjnymi ani logowania wielu użytkowników. Zakres został dobrany tak, aby program był możliwy do wykonania jako projekt studencki, ale jednocześnie miał rozbudowaną bazę danych i kilka powiązanych modułów.

\section{Cel projektu}

Celem projektu było przygotowanie aplikacji desktopowej w języku C\#, która umożliwia podstawową obsługę hotelu. System ma pozwalać na przeglądanie danych hotelowych, dodawanie gości, tworzenie rezerwacji, sprawdzanie dostępności pokoi, rejestrowanie płatności oraz zarządzanie usługami dodatkowymi.

\section{Zakres realizowanych funkcjonalności}

W projekcie zrealizowano podstawowy zakres funkcjonalności potrzebny do pracy recepcji hotelowej. Użytkownik może:

\begin{itemize}
  \item przeglądać dane w tabelach,
  \item dodawać nowych gości hotelowych,
  \item przeglądać pokoje oraz typy pokoi,
  \item tworzyć rezerwacje dla wybranego gościa i pokoju,
  \item sprawdzać dostępność pokoi w podanym terminie,
  \item rejestrować płatności dla wybranej rezerwacji,
  \item przeglądać listę pracowników,
  \item dodawać nowe usługi hotelowe,
  \item edytować istniejące usługi hotelowe,
  \item sprawdzać szybkie statystyki na panelu głównym,
  \item poruszać się po aplikacji za pomocą zakładek.
\end{itemize}

\section{Przegląd podobnych rozwiązań}

Pierwszym podobnym rozwiązaniem jest Booking.com. Jest to duża platforma internetowa do wyszukiwania i rezerwowania noclegów. W porównaniu z takim rozwiązaniem mój projekt nie obsługuje klientów przez Internet, tylko skupia się na lokalnej pracy recepcji.

Drugim przykładem jest system KWHotel. Jest to program przeznaczony do zarządzania obiektami noclegowymi, rezerwacjami i pokojami. Mój projekt realizuje podobny ogólny pomysł, ale w prostszym zakresie i bez funkcji biznesowych, takich jak faktury czy synchronizacja z portalami.

Trzecim przykładem jest Cloudbeds albo inny system PMS dla hoteli. Takie systemy są rozbudowane i często działają w chmurze. Mój projekt został wykonany jako lokalna aplikacja desktopowa, dlatego jest prostszy, ale zawiera najważniejsze elementy: pokoje, rezerwacje, gości, płatności i usługi.
